% block: machine-learning-heading page:68 type:section
\subsection{{Machine Learning}}\label{sec:machine-learning}
{Linear systems arise in the field of machine learning. In this section, through two examples: least-squares regression and least-squares classification, we will see the importance of linear systems.}
% block: least-squares-regression-heading page:68 type:subsection
\subsubsection{{Least-squares Regression}}\label{subsec:least-squares-regression}
{In this subsection, we consider one of the simplest supervised learning: least-squares learning. Given training data set $(\boldsymbol{x}_i,y_i)$ ($i=1,\ldots,n$) and a function $f_{\boldsymbol{\theta}}(\boldsymbol{x})$. Then the least-squares learning learns a parameter $\boldsymbol{\theta}\in\mathbb{R}^m$ such that the following error function is minimized:}
\begin{equation}\label{eq:least-squares-loss}
J(\boldsymbol{\theta})=\sum_{i=1}^n(y_i-f_{\boldsymbol{\theta}}(\boldsymbol{x}_i))^2.
\end{equation}
{If $f_{\boldsymbol{\theta}}(\boldsymbol{x})$ is given by the linear combination of functions $f_1(\boldsymbol{x}),\ldots,f_m(\boldsymbol{x})$ of the form}
\begin{equation}\label{eq:regression-basis-model}
f_{\boldsymbol{\theta}}(\boldsymbol{x})=\theta_1f_1(\boldsymbol{x})+\theta_2f_2(\boldsymbol{x})+\theta_mf_m(\boldsymbol{x}),
\end{equation}
{then the error function \eqref{eq:least-squares-loss} is rewritten as}\footnote{{In this section, we explicitly describe the 2-norm as $\|\cdot\|_2$. In machine learning, 1-norm $\|\cdot\|_1$ is also used, especially for obtaining a sparse solution of the least-squares problems of the form $\|\boldsymbol{y}-M\boldsymbol{\theta}\|_2^2+\lambda\|\boldsymbol{\theta}\|_1$.}}
\begin{equation}\label{eq:least-squares-matrix-loss}
J(\boldsymbol{\theta})=\|\boldsymbol{y}-M\boldsymbol{\theta}\|_2^2,
\end{equation}
{where $\boldsymbol{y}=[y_1,y_2,\ldots,y_n]^{\top}$, $\boldsymbol{\theta}=[\theta_1,\theta_2,\ldots,\theta_m]^{\top}$, and}
% block: regression-design-matrix page:69 type:ordinary-paragraph
\[M=\begin{bmatrix}
f_1(\boldsymbol{x}_1)&f_2(\boldsymbol{x}_1)&\cdots&f_m(\boldsymbol{x}_1)\\
f_1(\boldsymbol{x}_2)&f_2(\boldsymbol{x}_2)&\cdots&f_m(\boldsymbol{x}_2)\\
\vdots&\vdots&\ddots&\vdots\\
f_1(\boldsymbol{x}_n)&f_2(\boldsymbol{x}_n)&\cdots&f_m(\boldsymbol{x}_n)
\end{bmatrix}.\]
{For simplicity, column vectors of $M$ are assumed to be linearly independent. For minimizing the error function \eqref{eq:least-squares-matrix-loss}, the following decomposition is useful:}
\begin{equation}\label{eq:least-squares-square-completion}
\begin{aligned}
J(\boldsymbol{\theta})&=\|\boldsymbol{y}-M\boldsymbol{\theta}\|_2^2=(\boldsymbol{y}-M\boldsymbol{\theta})^{\top}(\boldsymbol{y}-M\boldsymbol{\theta})\\
&=\boldsymbol{v}_{\boldsymbol{\theta}}^{\top}(M^{\top}M)^{-1}\boldsymbol{v}_{\boldsymbol{\theta}}-(M^{\top}\boldsymbol{y})^{\top}(M^{\top}M)^{-1}(M^{\top}\boldsymbol{y})+\boldsymbol{y}^{\top}\boldsymbol{y}.
\end{aligned}
\end{equation}
{Here $\boldsymbol{v}_{\boldsymbol{\theta}}=M^{\top}\boldsymbol{y}-M^{\top}M\boldsymbol{\theta}$. Since the columns of $M$ are linearly independent, $M^{\top}M$ is a symmetric positive definite matrix, i.e., all the eigenvalues are positive, or equivalently $\boldsymbol{x}^{\top}M^{\top}M\boldsymbol{x}>0$ for all nonzero vectors $\boldsymbol{x}$. From this, $J(\boldsymbol{\theta})$ is minimized if $\boldsymbol{v}_{\boldsymbol{\theta}}=\boldsymbol{0}$, i.e., $M^{\top}\boldsymbol{y}-M^{\top}M\boldsymbol{\theta}=\boldsymbol{0}$. Thus we need to solve the following linear systems:}
\begin{equation}\label{eq:regression-normal-equations}
M^{\top}M\boldsymbol{\theta}=M^{\top}\boldsymbol{y}.
\end{equation}
% block: regression-example page:69 type:ordinary-paragraph
{As a numerical example, consider $y=(1/2)\sin(x)+\epsilon$, where $\epsilon$ is a noise. We generated a data set (training data set) with $(x_1,x_2,\ldots,x_{63})=(0,0.1,0.2,\ldots,6.2)$ and $(y_1,y_2,\ldots,y_{63})=(y(0),y(0.1),y(0.2),\ldots,y(6.2))$. The visualization of the data is shown in Fig.\ \ref{fig:regression-data}.}
{Next, we consider \eqref{eq:regression-basis-model}. For the basis functions $f_1(x),f_2(x),\ldots,f_m(x)$, we use}
\[f_k(x)=\begin{cases}\sin\frac12kx&(k:\text{even}),\\\cos\frac12kx&(k:\text{odd}).\end{cases}\]
{The unknown vector $\boldsymbol{\theta}$ of the model function \eqref{eq:regression-basis-model} is given by solving linear systems \eqref{eq:regression-normal-equations}, and the plots of $f_{\boldsymbol{\theta}}(x)$ with $m=5$ and $m=10$ are shown in Fig.\ \ref{fig:regression-fit}.}
% block: regression-data-figure page:69 type:figure
\begin{figure}[htbp]\centering\includegraphics[width=.48\linewidth]{figures/regression-data.png}
\caption{{The plot of the given data set}}\label{fig:regression-data}\end{figure}
% block: regression-fit-figure page:70 type:figure
\begin{figure}[htbp]\centering\includegraphics[width=.92\linewidth]{figures/regression-fit.png}
\caption{{The plots of $f_{\boldsymbol{\theta}}(x)$ in \eqref{eq:regression-basis-model} with $m=5$ and $m=10$. (a) The number of basis functions $m=5$. (b) The number of basis functions $m=10$.}}\label{fig:regression-fit}\end{figure}
% block: regression-overfitting page:70 type:ordinary-paragraph
{From Fig.\ \ref{fig:regression-fit}a, the function $f_{\boldsymbol{\theta}}(x)$ with $m=5$ is robust against the noise. On the other hand, the function with $m=10$ seems to try to fit the data with the given noise. This phenomenon is known as overfitting.}
{In what follows, one remedy for the overfitting is introduced, which also requires solving (parameterized) linear systems. We now consider the following constrained minimization problem of the form}
\begin{equation}\label{eq:constrained-regression}
\min_{\boldsymbol{\theta}}\|\boldsymbol{y}-M\boldsymbol{\theta}\|_2^2\quad\text{subject to }\|\boldsymbol{\theta}\|_2^2\le R.
\end{equation}
{If $R$ is infinite, the minimization problem is equivalent to \eqref{eq:least-squares-matrix-loss}. In order to solve \eqref{eq:constrained-regression}, the dual problem is useful. Let $f(\boldsymbol{\theta})=\|\boldsymbol{y}-M\boldsymbol{\theta}\|_2^2$ and $g(\boldsymbol{\theta})=\|\boldsymbol{\theta}\|_2^2-R$, then \eqref{eq:constrained-regression} can be rewritten as}
\begin{equation}\label{eq:abstract-constrained-regression}
\min_{\boldsymbol{\theta}}f(\boldsymbol{\theta})\quad\text{subject to }g(\boldsymbol{\theta})\le0.
\end{equation}
{Here we introduce the Lagrange function with $\lambda\ge0$:}
\[L(\boldsymbol{\theta},\lambda):=f(\boldsymbol{\theta})+\lambda g(\boldsymbol{\theta}).\]
{Since $\lambda\ge0$ and $g(\boldsymbol{\theta})\le0$, we obtain}
\[\max_{\lambda\ge0}L(\boldsymbol{\theta},\lambda)=\begin{cases}f(\boldsymbol{\theta})&\text{(if }g(\boldsymbol{\theta})\le0),\\+\infty&\text{(otherwise)}.\end{cases}\]
{Thus \eqref{eq:abstract-constrained-regression} is equivalent to}
\[\min_{\boldsymbol{\theta}}\max_{\lambda\ge0}L(\boldsymbol{\theta},\lambda).\]
% block: regression-duality page:71 type:ordinary-paragraph
{We now consider the following dual problem:}
\begin{equation}\label{eq:regression-dual-problem}
\max_{\lambda\ge0}\min_{\boldsymbol{\theta}}L(\boldsymbol{\theta},\lambda).
\end{equation}
{If $f$ and $g$ are convex, and certain conditions are satisfied, we have}
\[\max_{\lambda\ge0}\min_{\boldsymbol{\theta}}L(\boldsymbol{\theta},\lambda)=\min_{\boldsymbol{\theta}}\max_{\lambda\ge0}L(\boldsymbol{\theta},\lambda),\]
{which is called the strong duality. This is related to the Lagrangian duality theory; see \BookCite{26} for the details. The optimization problem \eqref{eq:regression-dual-problem} reads}
\begin{equation}\label{eq:regularized-regression-dual}
\max_{\lambda\ge0}\min_{\boldsymbol{\theta}}L(\boldsymbol{\theta},\lambda)\Leftrightarrow
\max_{\lambda\ge0}\min_{\boldsymbol{\theta}}\|\boldsymbol{y}-M\boldsymbol{\theta}\|_2^2+\lambda(\|\boldsymbol{\theta}\|_2^2-R).
\end{equation}
% block: regularized-square-completion page:71 type:ordinary-paragraph
{Similar to \eqref{eq:least-squares-square-completion}, the error function can be decomposed as}
\[\begin{aligned}
&\|\boldsymbol{y}-M\boldsymbol{\theta}\|_2^2+\lambda(\|\boldsymbol{\theta}\|_2^2-R)\\
&=\boldsymbol{w}_{\boldsymbol{\theta}}^{\top}(M^{\top}M+\lambda I)^{-1}\boldsymbol{w}_{\boldsymbol{\theta}}-(M^{\top}\boldsymbol{y})^{\top}(M^{\top}M+\lambda I)^{-1}M^{\top}\boldsymbol{y}+\boldsymbol{y}^{\top}\boldsymbol{y}-\lambda R,
\end{aligned}\]
{where $\boldsymbol{w}_{\boldsymbol{\theta}}=M^{\top}\boldsymbol{y}-(M^{\top}M+\lambda I)\boldsymbol{\theta}$. Since $\lambda\ge0$, matrix $M^{\top}M+\lambda I$ is symmetric positive definite. Thus $\boldsymbol{\theta}$ can be obtained by solving the following linear systems:}
\begin{equation}\label{eq:regularized-normal-equations}
(M^{\top}M+\lambda I)\boldsymbol{\theta}=M^{\top}\boldsymbol{y}.
\end{equation}
{Let $\boldsymbol{\theta}(\lambda)$ be the solution of the linear systems \eqref{eq:regularized-normal-equations}. Then \eqref{eq:regression-dual-problem} is equivalent to}
\[\max_{\lambda\ge0}L(\boldsymbol{\theta}(\lambda),\lambda).\]
{For approximately solving the above maximization problem, one may set $0\le\lambda_1<\lambda_2<\cdots<\lambda_m$ and then compute}
\[\max\{L(\boldsymbol{\theta}(\lambda_1),\lambda_1),L(\boldsymbol{\theta}(\lambda_2),\lambda_2),\ldots,L(\boldsymbol{\theta}(\lambda_m),\lambda_m)\}.\]
{In this case, we need $\boldsymbol{\theta}(\lambda_1),\boldsymbol{\theta}(\lambda_2),\ldots,\boldsymbol{\theta}(\lambda_m)$, i.e.,}
\begin{equation}\label{eq:regression-shifted-systems}
(M^{\top}M+\lambda_k I)\boldsymbol{\theta}(\lambda_k)=M^{\top}\boldsymbol{y}\quad k=1,2,\ldots,m.
\end{equation}
{Equations \eqref{eq:regression-shifted-systems} are known as shifted linear systems. Krylov subspace methods for solving shifted linear systems are described in Chap.\ \ref{chap:shifted-linear-systems}.}
% block: regularization-numerical-example pages:71,72 type:ordinary-paragraph
{Here, we show below a numerical example whose data is the same as in Fig.\ \ref{fig:regression-data}. Solving linear systems \eqref{eq:regularized-normal-equations} with a given $\lambda$ yields the unknown vector $\boldsymbol{\theta}$ of the model function \eqref{eq:regression-basis-model}. The plot of $f_{\boldsymbol{\theta}}(x)$ with $m=10$ and $\lambda=1.0$ is shown in Fig.\ \ref{fig:regularized-regression-fit}b. For reference and convenience, Fig.\ \ref{fig:regression-fit}b is shown again in Fig.\ \ref{fig:regularized-regression-fit}a.}
{In Fig.\ \ref{fig:regularized-regression-fit}b, we see the result is now robust against the noise and thus seems to avoid the overfitting that is seen in Fig.\ \ref{fig:regularized-regression-fit}a.}
% block: regularized-regression-fit-figure page:72 type:figure
\begin{figure}[htbp]\centering\includegraphics[width=.92\linewidth]{figures/regularized-regression-fit.png}
\caption{{The plots of $f_{\boldsymbol{\theta}}(x)$ in \eqref{eq:regression-basis-model} with $m=10$ and $m=10$ with $\lambda=1.0$. (a) $m=10$. (b) $m=10$, $\lambda=1.0$.}}\label{fig:regularized-regression-fit}\end{figure}
% block: probabilistic-classification-heading page:72 type:subsection
\subsubsection{{Least-squares Classification}}\label{subsec:least-squares-classification}
{In this subsection, least-squares probabilistic classification \BookCite{178} (see also \BookCite{148} for the summary) is described.}
{Let $\mathcal{X}\subset\mathbb{R}^d$ and $\mathcal{Y}=\{1,2,\ldots,c\}$, where $c$ is the number of classes. The problem of the classification is to classify $\boldsymbol{x}\in\mathcal{X}$ into classes $y\in\mathcal{Y}$, based on a given dataset (training dataset) $(\boldsymbol{x}_i,y_i)\in\mathcal{X}\times\mathcal{Y}$ for $i=1,2,\ldots,n$. Here $(\boldsymbol{x}_i,y_i)$ means that the data $\boldsymbol{x}_i\in\mathcal{X}$ belongs to the class $y_i\in\mathcal{Y}$.}
{The probabilistic pattern recognition is to estimate the class-posterior probability $p(y|\boldsymbol{x})$ from the training data $\{(\boldsymbol{x}_i,y_i)\}_{i=1}^n$. Once $p(y|\boldsymbol{x})$ is obtained, new data $\boldsymbol{x}$ can be classified by}
\[\widehat y=\operatorname*{arg\,max}_{y\in\{1,2,\ldots,c\}}p(y|\boldsymbol{x}).\]
{Let $p(\boldsymbol{x})$ be the marginal density of $\boldsymbol{x}$ with $p(\boldsymbol{x})>0$ for all $\boldsymbol{x}$. Then the class-posterior probability $p(y|\boldsymbol{x})$ is written as}
\begin{equation}\label{eq:class-posterior-probability}
p(y|\boldsymbol{x})=\frac{p(\boldsymbol{x},y)}{p(\boldsymbol{x})},
\end{equation}
{since $p(\boldsymbol{x},y)=p(y|\boldsymbol{x})p(\boldsymbol{x})$ and the assumption $p(\boldsymbol{x})>0$. The least-squares probabilistic classification uses the following linear model $q(y|\boldsymbol{x};\boldsymbol{\theta}^{(y)})$ for the class-posterior probability $p(y|\boldsymbol{x})$:}
% block: classification-model page:73 type:ordinary-paragraph
\begin{equation}\label{eq:classification-basis-model}
q(y|\boldsymbol{x};\boldsymbol{\theta}^{(y)})=\boldsymbol{\theta}^{(y)\top}\boldsymbol{f}(\boldsymbol{x})=\theta_1^{(y)}f_1(\boldsymbol{x})+\theta_2^{(y)}f_2(\boldsymbol{x})+\cdots+\theta_b^{(y)}f_b(\boldsymbol{x}),
\end{equation}
{where $\boldsymbol{\theta}^{(y)}=[\theta_1^{(y)},\theta_2^{(y)},\ldots,\theta_b^{(y)}]^{\top}$, and $f_i(\boldsymbol{x})$ ($i=1,\ldots,b$) are the given basis functions, e.g.,}
\begin{equation}\label{eq:classification-gaussian-kernel}
f_i(\boldsymbol{x})=K(\boldsymbol{x},\boldsymbol{x}_i)=\exp\left(-\frac{\|\boldsymbol{x}-\boldsymbol{x}_i\|_2^2}{2h^2}\right).
\end{equation}
{Here $K(\boldsymbol{x},\boldsymbol{x}_i)$ is known as the Gaussian kernel and $h$ is a scalar parameter that is set by users. Parameter $\boldsymbol{\theta}^{(y)}$ is determined by approximately minimizing the following least squares error:}
\[J_y(\boldsymbol{\theta}^{(y)})=\frac12\int\left(q(y|\boldsymbol{x};\boldsymbol{\theta}^{(y)})-p(y|\boldsymbol{x})\right)^2p(\boldsymbol{x})\,\mathrm d\boldsymbol{x}.\]
% block: classification-loss-expansion page:73 type:ordinary-paragraph
{$J_y(\boldsymbol{\theta}^{(y)})$ can be rewritten as}
\begin{equation}\label{eq:classification-integral-loss}
\begin{aligned}
J_y(\boldsymbol{\theta}^{(y)})&=\frac12\int q(y|\boldsymbol{x};\boldsymbol{\theta}^{(y)})^2p(\boldsymbol{x})\,\mathrm d\boldsymbol{x}-\int q(y|\boldsymbol{x};\boldsymbol{\theta}^{(y)})p(y|\boldsymbol{x})p(\boldsymbol{x})\,\mathrm d\boldsymbol{x}+c\\
&=\frac12\int q(y|\boldsymbol{x};\boldsymbol{\theta}^{(y)})^2p(\boldsymbol{x})\,\mathrm d\boldsymbol{x}-\int q(y|\boldsymbol{x};\boldsymbol{\theta}^{(y)})p(\boldsymbol{x},y)\,\mathrm d\boldsymbol{x}+c\\
&=\frac12\int q(y|\boldsymbol{x};\boldsymbol{\theta}^{(y)})^2p(\boldsymbol{x})\,\mathrm d\boldsymbol{x}-\int q(y|\boldsymbol{x};\boldsymbol{\theta}^{(y)})p(y)p(\boldsymbol{x}|y)\,\mathrm d\boldsymbol{x}+c,
\end{aligned}
\end{equation}
{where $c=(1/2)\int p(y|\boldsymbol{x})^2p(\boldsymbol{x})\,\mathrm d\boldsymbol{x}$, and \eqref{eq:class-posterior-probability} was used in the second equation. By using the average of the training data $(\boldsymbol{x}_1,y_1),(\boldsymbol{x}_2,y_2),\ldots,(\boldsymbol{x}_n,y_n)$, we have the following approximation:}
\[\begin{aligned}
\int q(y|\boldsymbol{x};\boldsymbol{\theta}^{(y)})^2p(\boldsymbol{x})\,\mathrm d\boldsymbol{x}&\approx\frac1n\sum_{i=1}^nq(y|\boldsymbol{x}_i;\boldsymbol{\theta}^{(y)})^2,\\
\int q(y|\boldsymbol{x};\boldsymbol{\theta}^{(y)})p(y)p(\boldsymbol{x}|y)\,\mathrm d\boldsymbol{x}&\approx\frac1{n_y}\sum_{i:y_i=y}q(y|\boldsymbol{x}_i;\boldsymbol{\theta}^{(y)})p(y)\approx\frac1n\sum_{i:y_i=y}q(y|\boldsymbol{x}_i;\boldsymbol{\theta}^{(y)}).
\end{aligned}\]
{In the last approximation, we used $p(y)\approx n_y/n$, where $n_y$ is the number of training data that belong to the class $y\in\mathcal{Y}$. Note that the symbol $\sum_{i:y_i=y}$ means the sum over $i$ such that $y_i=y$ for the given class $y\in\mathcal{Y}$. Using this approximation to \eqref{eq:classification-integral-loss} together with regularization term $(\lambda/2n)\|\boldsymbol{\theta}^{(y)}\|_2^2$, see also \eqref{eq:regularized-regression-dual}, we have the following new minimization problem:}
\begin{equation}\label{eq:classification-empirical-loss}
\widehat J_y(\boldsymbol{\theta}^{(y)})=\frac1{2n}\sum_{i=1}^nq(y|\boldsymbol{x}_i;\boldsymbol{\theta}^{(y)})^2-\frac1n\sum_{i:y_i=y}q(y|\boldsymbol{x}_i;\boldsymbol{\theta}^{(y)})+\frac\lambda{2n}\|\boldsymbol{\theta}^{(y)}\|_2^2,
\end{equation}
% block: classification-matrix-loss page:74 type:ordinary-paragraph
{where the constant $c$ was omitted because it does not affect the minimization problem. Let $(M)_{i,j}=f_i(\boldsymbol{x}_j)$. Then from \eqref{eq:classification-basis-model}, approximate error function \eqref{eq:classification-empirical-loss} can be rewritten as}
\begin{equation}\label{eq:classification-matrix-loss}
\widehat J_y(\boldsymbol{\theta}^{(y)})=\frac1{2n}\boldsymbol{\theta}^{(y)\top}M^{\top}M\boldsymbol{\theta}^{(y)}-\frac1n\boldsymbol{\theta}^{(y)\top}M^{\top}\boldsymbol{p}^{(y)}+\frac\lambda{2n}\boldsymbol{\theta}^{(y)\top}\boldsymbol{\theta}^{(y)},
\end{equation}
{where $\boldsymbol{p}^{(y)}$ is a vector whose $i$th element $p_i^{(y)}$ is defined by}
\[p_i^{(y)}=\begin{cases}1&\text{(if }y_i=y),\\0&\text{(if }y_i\ne y).\end{cases}\]
{Similar to the previous subsection, \eqref{eq:classification-matrix-loss} can be decomposed as}
\[\begin{aligned}
\widehat J_y(\boldsymbol{\theta}^{(y)})&=\frac1{2n}\left(M^{\top}\boldsymbol{p}^{(y)}-G\boldsymbol{\theta}^{(y)}\right)^{\top}G^{-1}\left(M^{\top}\boldsymbol{p}^{(y)}-G\boldsymbol{\theta}^{(y)}\right)\\
&\quad-\frac1{2n}\boldsymbol{p}^{(y)\top}MG^{-1}M^{\top}\boldsymbol{p}^{(y)},
\end{aligned}\]
{where $G=M^{\top}M+\lambda I$. If columns of $M$ are linearly independent, then $G$ is symmetric positive definite for all $\lambda\ge0$. Thus $\widehat J_y(\boldsymbol{\theta}^{(y)})$ is minimized when $M^{\top}\boldsymbol{p}^{(y)}-G\boldsymbol{\theta}^{(y)}=\boldsymbol{0}$, which leads to}
\begin{equation}\label{eq:classification-normal-equations}
(M^{\top}M+\lambda I)\boldsymbol{\theta}^{(y)}=M^{\top}\boldsymbol{p}^{(y)}.
\end{equation}
% block: classification-normalization page:74 type:ordinary-paragraph
{After obtaining the solution of $\boldsymbol{\theta}^{(y)}$, we have \eqref{eq:classification-basis-model}. Then, the following normalization gives the approximated class probability:}
\[\widehat p(y|\boldsymbol{x})=\frac{\max(0,q(y|\boldsymbol{x};\boldsymbol{\theta}^{(y)}))}{\sum_{\widetilde y=1}^c\max(0,q(\widetilde y|\boldsymbol{x};\boldsymbol{\theta}^{(\widetilde y)}))}.\]
{In what follows, a simple numerical example is illustrated to see how the least-squares probabilistic classification works. We use the training data as given in Fig.\ \ref{fig:classification-result}a, where there are two classes: $\circ$ denotes class $y=1$ and $\times$ denotes class $y=2$. From \eqref{eq:classification-basis-model} and \eqref{eq:classification-gaussian-kernel}, the following model is used:}
\[q(y|\boldsymbol{x};\boldsymbol{\theta}^{(y)})=\sum_{i:y_i=y}\theta_i^{(y)}K(\boldsymbol{x},\boldsymbol{x}_i).\]
{This means that we will obtain two function $q(y|\boldsymbol{x};\boldsymbol{\theta}^{(y)})$ for $y=1$ and $y=2$, after solving two linear systems \eqref{eq:classification-normal-equations} for $y=1$ and $y=2$. The result of the least-squares probabilistic classification is given in Fig.\ \ref{fig:classification-result}, where the number of samples $n=60$, and the two parameters $\lambda$ in \eqref{eq:classification-empirical-loss} and $h$ in \eqref{eq:classification-gaussian-kernel} are $\lambda=0.1$ and $h=1$.}
% block: classification-result-figure page:75 type:figure
\begin{figure}[htbp]\centering\includegraphics[width=.92\linewidth]{figures/classification-result.png}
\caption{{The result of the least-squares probabilistic classification. (a) Training samples of two classes. (b) The class-posterior probabilities.}}\label{fig:classification-result}\end{figure}
% block: classification-result-discussion page:75 type:ordinary-paragraph
{Fig.\ \ref{fig:classification-result}a shows the distribution of two classes of samples as mentioned before, and Fig.\ \ref{fig:classification-result}b is the corresponding result. From Fig.\ \ref{fig:classification-result}b, we see that the probability of the point $x=0$ that belongs to class 1 ($y=1$) is about 50\%, i.e., $p(y=1|x=0)\approx0.5$.}
